\chapter{Instalación y Primeros Pasos}
\label{chap:installation}

{\huge$\bigstar$} Principiante

\section{Requisitos del Sistema}

LisPy está diseñado para Windows x64. Los requisitos mínimos son:

\begin{itemize}
\item Sistema operativo: Windows 10 o superior (x64)
\item Opcional: CPython 3.14 instalado (para soporte completo del puente Python)
\item Espacio en disco: ~50 MB para la instalación base
\end{itemize}

\section{Descarga e Instalación}

Descarga el binario \texttt{lispy.exe} desde el repositorio oficial y
colócalo en un directorio de tu elección. Opcionalmente, añade dicho
directorio a la variable de entorno \texttt{PATH} para acceso global.

\section{Verificación}

Abre una terminal y ejecuta:

\begin{lispcode}
lispy.exe --version
\end{lispcode}

Deberías ver información sobre la versión y configuración del intérprete.

\section{Entornos Virtuales Python}

Si deseas aislar las dependencias Python de tu proyecto:

\begin{lstlisting}[style=python,title=Terminal]
python -m venv mi_entorno
lispy.exe --venv mi_entorno
\end{lstlisting}

\section{Primer Comando}

Inicia el REPL:

\begin{lstlisting}[style=python,title=Terminal]
lispy.exe --repl
\end{lstlisting}

En el prompt \texttt{LisPy>}, escribe:

\begin{lispcode}
(+ 1 2)
\end{lispcode}

LisPy responderá con \texttt{3}.

\section{El REPL Interactivo}

El REPL de LisPy ofrece los siguientes comandos:

\begin{itemize}
\item \texttt{:help} -- Muestra la ayuda del REPL
\item \texttt{:quit} -- Sale del REPL
\item \texttt{:pip} -- Gestión de paquetes Python
\item \texttt{:trace} -- Activa/desactiva traza de ejecución
\end{itemize}

\section{Ejemplos}

\begin{examplebox}
\begin{lispcode}
;; hello.lp - Primer programa
(format t "Hello, LisPy!~%")
\end{lispcode}
Salida: \texttt{Hello, LisPy!}
\end{examplebox}

\begin{examplebox}
\begin{lispcode}
;; Uso de --venv con un entorno virtual preexistente
;; Terminal:
;;   lispy.exe --venv mi_entorno script.lp

;; Instalación de numpy desde el REPL
;; LisPy> :pip install numpy
\end{lispcode}
\end{examplebox}

